\chapter{FPGA 代码设计与实现}
\label{chap:hardware}

\tikzset{
  fpblock/.style={draw=TemplateBlue,rounded corners=2pt,fill=TemplateGray,
    align=center,font=\fontsize{9.5}{12.5}\selectfont,minimum height=10mm,text width=4cm,inner sep=4pt},
  fpstate/.style={fpblock,text width=3.1cm},
  fparrow/.style={-{Stealth[length=2mm]},thick,draw=TemplateBlue},
  fpnote/.style={align=center,font=\fontsize{8.5}{11}\selectfont,text=TemplateBlue}
}

\section{设计范围与 RTL 总体结构}
本章以 \path{../Pgl50h/EmbFPGA.pds} 登记的工程及
\path{src/top/hdmi_ddr3_top.v} 默认参数为依据，说明第三章算法在
PGL50H-FBG484-6 上的 RTL 实现。系统输入、输出均为
$1920\times1080$、60 Hz 的 HDMI 视频，外部并行接口采用 RGB888，
DDR3 内的源帧和结果帧采用 RGB565。算法按完整源帧调度，显示端按独立光栅连续扫描，
两者通过结果帧槽解耦；60 Hz 输出时序本身不代表每秒产生 60 个不同的矫正结果。

默认参数为 \path{UNDISTORT_ENABLE=1}、\path{POSE_CORRECTION_ENABLE=1}、
\path{POSE_TILED_ENABLE=1}、\path{INTERPOLATION_MODE=1}，
即启用双槽分块逆投影与最近邻采样。输入预处理选择
\path{ADAPTIVE_PREPROCESS_MODE=3}，启用帧自适应 Gamma；
\path{PC_CONTROL_ENABLE=1} 支持运行时控制。
\path{POSE_RECIPROCAL_LUT_ENABLE=0} 选择恢复除法流水线，
\path{OUTPUT_BYPASS=0} 选择处理后输出。
双线性插值由保留的串行路径实现，不属于本章默认分块数据通路。

图~\ref{fig:fpga-architecture} 给出按实际实例关系整理的 RTL 结构原理图。
图中的 DDR3 控制器以及时钟通路使用的 PLL 采用厂商 IP；
像素处理、帧管理和仲裁功能由工程 RTL 实现。
\begin{figure}[htbp]
\centering
\begin{tikzpicture}
\node[fpblock] (rx) at (0,0) {MS7200 并行输入\\RGB888、DE、VS};
\node[fpblock] (gamma) at (5,0) {输入寄存 + 自动 Gamma\\\texttt{pixclk\_in}};
\node[fpblock] (capture) at (10,0) {RGB565 / 完整帧采集\\打包 → 采集异步 FIFO};
\node[fpblock] (render) at (0,-2.3) {双槽分块渲染器\\联合逆投影 / 源缓存 / 采样};
\node[fpblock] (arb) at (5,-2.3) {DDR 读写仲裁器\\采集/渲染写，显示/渲染读};
\node[fpblock] (ddr) at (10,-2.3) {\texttt{DDR3\_50H} IP\\用户数据 256 bit / DDR3};
\node[fpblock] (display) at (0,-4.6) {结果帧预取 / 显示 FIFO\\RGB565 解包};
\node[fpblock] (output) at (5,-4.6) {RGB888 展开 / 输出寄存\\独立 1080p60 光栅};
\node[fpblock] (tx) at (10,-4.6) {MS7210 并行输出\\RGB888、DE、HS、VS};
\draw[fparrow] (rx)--(gamma);
\draw[fparrow] (gamma)--(capture);
\draw[fparrow] (capture.south)--++(0,-0.65)-| (arb.north);
\draw[fparrow] (arb.east)--(ddr.west);
\draw[fparrow] (ddr.south)--++(0,-0.45)-| (arb.south east);
\draw[fparrow,<->] (render.east)--(arb.west);
\draw[fparrow] (arb.south)--++(0,-0.65)-| (display.north);
\draw[fparrow] (display)--(output);
\draw[fparrow] (output)--(tx);
\node[fpnote] at (0,-1.2) {帧级锁存 IMU 姿态与几何开关};
\node[fpnote] at (10,0.95) {PC 掩码在输入 SOF 应用};
\node[fpnote] at (5,-5.6) {帧槽管理、采集、渲染和显示预取由 \texttt{fram\_buf\_pose} 组织};
\end{tikzpicture}
\caption{默认配置的 RTL 数据通路结构原理图（按源码连接重绘）}
\label{fig:fpga-architecture}
\end{figure}


\begin{longtable}{p{6.0cm}p{9.0cm}}
\caption{默认数据通路中的主要模块及职责}\label{tab:modules}\\
\toprule 模块（相对 \path{src/}） & 实现职责 \\\midrule\endfirsthead
\toprule 模块（相对 \path{src/}） & 实现职责 \\\midrule\endhead
\bottomrule\endfoot
\path{top/hdmi_ddr3_top.v} & 连接板级接口、PLL、DDR3 IP、视频处理和控制通路；按参数选择生成分支。\\
\path{board/ms72xx_ctl.v} & 组织 MS7200、MS7210 的寄存器初始化及两路 I\textsuperscript{2}C 接口。\\
\path{video/preprocess/adaptive_gamma_stream.v} & 原始帧亮度统计、目标曲线选择、三通道同步 LUT 查询及旁路对齐。\\
\path{video/format/rgb888_to_rgb565.v} & 对齐 VS、DE 与 RGB，寄存输出 RGB565 写接口。\\
\path{video/pose/fram_buf_pose.v} & 完整源帧采集、帧槽分配、姿态快照、渲染启动、结果发布与显示读出。\\
\path{video/pose/pose_renderer_tiled.v} & 双块槽调度、逐像素投影、源区域预取、最近邻采样及结果突发写回。\\
\path{video/pose/pose_inverse_project_q14.v} & 帧姿态系数准备、定点逆投影、镜头模型及坐标有效性判定。\\
\path{video/pose/pose_ddr_arbiter.v} & 两个写客户端、两个读客户端共享 DDR3 用户接口；核验长度、ID 和 LAST。\\
\path{video/pose/pose_async_fifo.v} & 采集和显示数据的双时钟缓冲，含同步 RAM 预取与 Gray 指针同步。\\
\path{video/timing/sync_vg.v}、\path{video_raster_delay.v} & 独立生成 1080p60 光栅，延迟显示 DE；二者位于同一目录。\\
\path{video/format/rgb565_to_rgb888.v}、\path{video/preprocess/gamma_correct_rgb888.v} & RGB 展开及输出寄存；自动 Gamma 模式关闭后级固定 Gamma 运算。\\
\path{imu/imu_protocol_50m.v}、\path{imu/attitude_cdc.v} & IMU 串口解析及姿态快照跨域。\\
\path{control/cp2102_control.v}、\path{control/video_control_cdc.v} & PC 协议解析、串口应答及输入帧边界应用控制掩码。\\
\path{control/gamma_telemetry.v}、\path{control/cdc_snapshot.v} & 形成同帧统计快照，握手跨域后发送遥测。\\
\end{longtable}
工程还保留线性帧缓存、单独去畸变及串行姿态矫正等生成分支。
独立文件 \path{pose_frame_bank_allocator.v} 和 \path{pose_tile_read_scheduler.v}
虽登记在工程中，默认路径未实例化；本章的帧槽、块读取调度分别对应
\path{fram_buf_pose.v} 和 \path{pose_renderer_tiled.v} 内部逻辑。
\sourcefile{\path{../Pgl50h/EmbFPGA.pds}；\path{../Pgl50h/src/top/hdmi_ddr3_top.v}}

\section{顶层接口、时钟与启动控制}
\subsection{顶层端口及接口边界}
表~\ref{tab:fpga-top-ports} 中方向均相对于 FPGA。
视频接口使用 DE 指示有效像素，输入 HS 不参与默认路径的逐行地址计数；
完整帧内有效像素按 DE 顺序打包，不进行水平抽列或缩放。
\begin{table}[htbp]\centering\small
\caption{顶层主要物理接口}\label{tab:fpga-top-ports}
\begin{tabularx}{\textwidth}{p{5.3cm}p{1.2cm}Y}
\toprule 端口组及位宽 & 方向 & 含义 \\\midrule
\path{sys_clk}、\path{video_ref_clk} & 输入 & 板载 50 MHz 系统时钟、27 MHz 视频参考时钟。\\
\path{pixclk_in}、\path{vs_in/hs_in/de_in}、\path{r_in/g_in/b_in[7:0]} & 输入 & MS7200 输出的源同步像素总线，1080p60 下约 148.5 MHz。\\
\path{pixclk_out}、\path{vs_out/hs_out/de_out}、\path{r_out/g_out/b_out[7:0]} & 输出 & 驱动 MS7210 的并行视频总线。\\
\path{imu_uart_rx} & 输入 & IMU UART，921600 bit/s、8N1。\\
\path{pc_uart_rx}、\path{pc_uart_tx} & 入/出 & 板载 CP2102 UART，115200 bit/s、8N1。\\
\path{sys_led[7:0]} & 输出 & 低三位依次显示镜头、姿态、Gamma 的已应用使能，其余位为 0。\\
\path{rstn_out}、两路 \path{iic*_scl/sda} & 出/双向 & 桥芯片复位与寄存器配置，SDA 通过三态接口连接。\\
\path{mem_a[14:0]}、\path{mem_ba[2:0]}、命令/时钟端口 & 输出 & DDR3 行列/Bank 地址、差分时钟及命令控制。\\
\path{mem_dq[31:0]}、\path{mem_dqs(_n)[3:0]}、\path{mem_dm[3:0]} & 双向/出 & 32 bit DDR3 物理数据、四组差分数据选通与字节屏蔽。\\
\bottomrule
\end{tabularx}
\end{table}

\subsection{时钟域划分与复位顺序}
系统分为四类工作域，如表~\ref{tab:fpga-clocks} 所示。
输入、输出像素时钟即使标称频率相同，也来自不同信号源，不能直接混用。
\begin{table}[htbp]\centering\small
\caption{实际工作时钟与承载逻辑}\label{tab:fpga-clocks}
\begin{tabularx}{\textwidth}{p{3.0cm}p{3.0cm}Y}
\toprule 时钟域 & 来源与频率 & 承载逻辑 \\\midrule
\path{sys_clk} & 板载 50 MHz & UART 接收/发送、IMU 协议、PC 控制。\\
\path{cfg_clk} & $27\times44/119\approx9.983$ MHz & 启动等待及 MS72xx I\textsuperscript{2}C 配置。\\
\path{pixclk_in} & HDMI 接收像素时钟 & 输入寄存、自动 Gamma、RGB565 打包、源帧姿态锁存。\\
\path{ddr_core_clk} & DDR3 IP 用户时钟 100 MHz & 采集写入、帧槽管理、投影/渲染、DDR 仲裁及显示预取。\\
\path{pixel_clk} & $27\times44/8=148.5$ MHz & 输出光栅、RGB565 解包、输出流水线。\\
\bottomrule
\end{tabularx}
\end{table}
视频 PLL 锁定后，\path{cfg_clk} 域计数至 10000，约等待 1 ms，
再拉高 \path{rstn_out}。\path{ms72xx_ctl} 先完成接收端 MS7200 的配置，
再以接收端完成信号释放 MS7210 配置状态机；两路 I\textsuperscript{2}C
使用独立驱动，按约 400 kHz 参数配置，寄存器地址低字节先发送。
DDR3 IP 同时以 \path{rstn_out} 释放复位并完成校准。

系统、输入像素和输出像素逻辑设置各自的两级复位释放同步寄存器；
帧缓存与异步 FIFO 以 \path{ddr_init_done} 为复位条件，分别在像素域释放。
跨域数据通路只有在相应复位释放且有效标志成立后使用。
输出光栅由 \path{video_reset_n} 控制，独立于 DDR 校准和首帧完成状态运行，
因此初始化尚未完成时仍能提供标准行场时序，并用诊断颜色表示状态。

\section{视频输入、自适应 Gamma 与输出时序}
\subsection{固定延迟的输入数据通路}
默认输入通路如图~\ref{fig:fpga-gamma}。
原始 RGB、DE 和低有效 VS 首先在 \path{pixclk_in} 上升沿寄存；
送入自动 Gamma 时将 VS 取反，以高有效 VS 的上升沿表示帧起始 SOF。
Gamma 的两个寄存级同时处理 RGB、完整 VS 电平和 DE，
再由 \path{rgb888_to_rgb565} 寄存成写接口。总计四拍输入延迟，
其中格式转换为
\begin{equation}
P_{565}=\{R[7:3],\ G[7:2],\ B[7:3]\}.
\end{equation}
格式转换输出 \path{wr_en}、\path{wr_data[15:0]} 和 \path{vs_out}
进入帧缓存；SOF 在帧缓存入口恢复为高有效 \path{wr_fsync}。
Gamma 关闭时使用同延迟旁路，保证开关前后帧同步与像素位置不变。
\begin{figure}[htbp]
\centering
\begin{tikzpicture}
\node[fpblock] (in) at (0,0) {输入 I/O 寄存\quad 1 拍\\RGB888、DE、VS};
\node[fpblock] (lut) at (5,0) {Gamma 寄存 / 同步 LUT\\RGB、DE、VS 共 2 拍};
\node[fpblock] (convert) at (10,0) {RGB888 → RGB565\quad 1 拍\\\texttt{wr\_data / wr\_en / VS}};
\node[fpblock] (stats) at (0,-2.3) {原始亮度与 DE 统计\\$S,N,N_{\rm dark},N_{\rm high}$};
\node[fpblock] (target) at (5,-2.3) {移位 / 加法阈值比较\\连续两帧相同 → 目标曲线};
\node[fpblock] (active) at (10,-2.3) {SOF 锁存使能 / 活动曲线\\向旧目标每帧移动至多一级};
\draw[fparrow] (in)--(lut);
\draw[fparrow] (lut)--(convert);
\draw[fparrow] (in)--(stats);
\draw[fparrow] (stats)--(target);
\draw[fparrow] (target)--(active);
\draw[fparrow,dashed] (active.north)--(10,-1.1)-| (lut.south);
\node[fpnote] at (5,1.0) {关闭 Gamma 选择同延迟旁路；总输入延迟为 4 拍，启动间隔为 1 拍};
\node[fpnote] at (5,-3.35) {统计在下一 SOF 发布，新分类只能影响后续完整帧；首帧使用恒等曲线};
\end{tikzpicture}
\caption{自适应 Gamma 的像素流水线与帧级控制回路}
\label{fig:fpga-gamma}
\end{figure}


\subsection{Gamma 的统计、控制与 LUT 电路}
\path{adaptive_gamma_stream} 的数据接口为
\path{in_r/g/b[7:0]}、\path{in_de/in_vs} 输入和对应
\path{out_r/g/b[7:0]}、\path{out_de/out_vs} 输出，流水线启动间隔为一拍，
接口不支持输入回压。控制端 \path{enable} 在 SOF 锁存，
输出 \path{active_curve[2:0]}、\path{target_curve[2:0]} 表示曲线状态。

亮度以 $Y=(R+2G+B)\gg2$ 计算，仅对原始输入的 DE 像素累加。
内部亮度和为 30 bit，像素数、暗像素数和高亮像素数为 22 bit，
分别覆盖 1080p 的最大亮度和 528768000 与像素数 2073600。
统计结果在下一帧 SOF 以 \path{stats_valid} 发布，四项统计扩展为 32 bit；
\path{stats_curve/stats_target/stats_enabled} 描述刚完成的统计帧。

曲线选择电路通过亮度和 $S$ 与常数倍像素数 $N$ 比较，避免实时求平均值除法。
在不满足混合光照判据时，按表~\ref{tab:fpga-gamma-select} 选曲线。
若暗像素 $Y<48$ 和高亮像素 $Y\ge208$ 均至少占 $1/8$，优先选择曲线 2。
\begin{table}[htbp]\centering\small
\caption{自动 Gamma 的分类规则及预存指数}\label{tab:fpga-gamma-select}
\begin{tabular}{ccc}
\toprule 分类条件（依次判定） & 曲线编号 & LUT 指数 $\gamma$ \\\midrule
$S<48N$ & 0 & 0.50\\
$48N\le S<80N$ & 1 & 0.65\\
$80N\le S<112N$ & 2 & 0.80\\
$S>208N$ & 6 & 1.90\\
$176N<S\le208N$ & 5 & 1.55\\
$144N<S\le176N$ & 4 & 1.25\\
其余 & 3 & 1.00\\
\bottomrule
\end{tabular}
\end{table}
连续两次分类相同才更新目标曲线；活动曲线仅在 SOF 向当时已有目标移动一级。
当前 SOF 发布的统计稍后才形成新目标，因此不会改变正在开始的这一帧。
复位时活动与目标均为恒等曲线 3。
\path{adaptive_gamma_lut} 按
$\operatorname{round}[255(v/255)^\gamma]$ 预存映射，
地址为曲线号与 8 bit 通道值拼接；每通道实例化一份 $2048\times8$ bit 同步 ROM。
关闭增强仍保留原始统计与曲线状态演进，像素选择延迟匹配的旁路值。
\sourcefile{\path{../Pgl50h/src/video/preprocess/adaptive_gamma_stream.v}；\path{adaptive_gamma_lut.v}}

\subsection{显示读出及行场对齐}
\path{sync_vg} 的行总长为 2200 像素，场总长为 1125 行，
有效区域为 $1920\times1080$；水平前肩/同步/后肩分别为 88/44/148，
垂直分别为 4/5/36。在 148.5 MHz 下，帧率为
$148.5\times10^6/(2200\times1125)=60$ Hz。
默认 \path{alg_rd_en=timing_de}，每个有效扫描位置都请求一个像素。

显示 FIFO 的 256 bit 数据按低位到高位拆成 16 个 RGB565 像素，
\path{fram_buf_pose} 内两拍寄存输出像素。
\path{video_raster_delay} 将光栅 DE 延迟同样的两拍，
\path{rgb565_to_rgb888} 用位复制组合展开到 RGB888，
后级 \path{gamma_correct_rgb888} 即使关闭运算仍保留一拍寄存，
最终顶层输出再寄存一拍。因此默认数据、DE、HS、VS 相对于输出时序源均延迟四拍。
自动 Gamma 生效时后级固定 Gamma 被旁路，避免重复校正。

最终视频寄存器放置在输出 I/O 单元；物理 \path{pixclk_out} 保持正相，
RGB/控制在上升沿更新，MS7210 按下降沿采样。
DE 无效时 RGB 清零；有效期内依优先级输出初始化诊断颜色：
发送端未完成为品红、DDR PLL 未锁定为蓝、DLL 未锁定为青、
DDR 未完成校准为红、尚无完整结果帧为黄，均就绪后输出处理数据。
\sourcefile{\path{../Pgl50h/src/top/hdmi_ddr3_top.v}；\path{../Pgl50h/src/video/timing/}}

\section{完整帧缓存与帧槽所有权}
\subsection{帧缓存接口和 DDR 地址组织}
\path{fram_buf_pose} 将采集、渲染与显示封装为统一帧缓存。
输入域接收 \path{wr_fsync/wr_en/wr_data[15:0]} 和已跨域的
\path{imu_yaw/imu_pitch[15:0]}、\path{lens_enable}；
输出域接收 \path{rd_fsync/rd_en}，返回
\path{vout_de/vout_data[15:0]}；DDR 域输出下节定义的用户命令。
\path{init_done} 在首个结果帧交给显示后置位，
\path{captured_frames/rendered_frames/displayed_frames} 分别计数采集提交、渲染完成和结果发布次数。
这些诊断计数及 \path{protocol_error} 等端口在当前顶层留空，未直接接入 UART 遥测。

一帧的有效数据量为
\begin{equation}
S_f=1920\times1080\times2=4147200\ \text{bytes},\qquad
N_b=\frac{1920\times1080}{16}=129600\ \text{beat}.
\end{equation}
DDR 用户地址以 32 bit word 为单位，一拍 256 bit 占 8 个地址单位，
同一帧相邻行首地址相差 $1920/2=960$。
帧槽 $b$ 的基址及其中像素 $(x,y)$ 所在 beat 地址为
\begin{equation}
A_b=b\,2^{22},\qquad
A(b,x,y)=A_b+960y+8\left\lfloor\frac{x}{16}\right\rfloor.
\label{eq:fpga-ddr-address}
\end{equation}
像素在 beat 内的 lane 为 $x\bmod16$，首像素位于数据低 16 bit。
基址转换为字节地址后如表~\ref{tab:fpga-banks}。
\begin{table}[htbp]\centering\small
\caption{DDR3 帧槽布局}\label{tab:fpga-banks}
\begin{tabular}{cccc}
\toprule 槽号 & 用户地址基址（word） & 字节地址基址 & 用途 \\\midrule
0 & \texttt{0x0000000} & \texttt{0x00000000} & 源帧 A\\
1 & \texttt{0x0400000} & \texttt{0x01000000} & 结果帧 A\\
2 & \texttt{0x0800000} & \texttt{0x02000000} & 结果帧 B\\
3 & \texttt{0x0C00000} & \texttt{0x03000000} & 源帧 B\\
\bottomrule
\end{tabular}
\end{table}
帧槽间隔为 16 MiB，仅有效帧载荷参与传输；DDR 帧槽与后文片内两个块槽是不同层次的缓存。

\subsection{源帧采集、提交与并行调度}
DDR 域翻转 \path{capture_request} 请求一帧，输入域同步该翻转，
在下一个 SOF 才开始采集，并同时锁存 yaw、pitch 和镜头使能。
每收齐 16 个有效像素产生一个 256 bit FIFO 写脉冲。
DDR 域从采集 FIFO 取数，先装满最多 16 beat 的客户端突发 RAM，再提交写命令。
未凑满 16 beat 的尾部在确认采集结束且 FIFO 排空后提交。

输入域收齐 2073600 个像素后发布完成翻转；DDR 域同步完成与失败标志，
再等待 8 个 DDR 时钟，使完成消息和最后一笔 FIFO 写指针均有时间到达。
只有 FIFO 无有效数据、占用计数为零、无未完成采集写、突发缓冲为空，
且已完成写入恰为 129600 beat 时，才接受该源帧。
输入提前出现新 SOF、FIFO 满导致写脉冲未被接受或帧长度不符，均使该帧被丢弃并重新请求。
HDMI 输入不能回压，FIFO 的 \path{wr_ready} 只用于检测接收失败。

默认快速路径的控制状态为
\path{C_REQUEST}（请求采集）、\path{C_CAPTURE}（收集并提交）和
\path{C_RENDER}（完整源帧等待调度）。后者需等待渲染器空闲、无待发布结果且无渲染错误，
再把源帧槽、姿态和镜头使能交给渲染器。
启动后采集槽在 0/3 间切换，下一帧采集与当前帧渲染重叠；
若新源帧已就绪而渲染器仍忙，就保持该源帧，并跳过后续输入帧，直到它可以调度。
\begin{figure}[htbp]
\centering
\begin{tikzpicture}
\node[fpblock] (cap) at (0,0) {源槽采集（0/3 轮换）\\SOF 锁存姿态并接收像素};
\node[fpblock] (ready) at (5,0) {完整源帧等待调度\\FIFO 排空、全部采集写结束};
\node[fpblock] (render) at (10,0) {源槽只读供渲染\\整帧像素与姿态保持不变};
\draw[fparrow] (cap)--(ready);
\draw[fparrow] (ready)--(render);
\draw[fparrow] (render.north)--++(0,0.8)-| (cap.north);
\node[fpnote,fill=white] at (5,1.0) {前一渲染完成后源槽可再次使用；另一源槽可并行采集};
\node[fpblock] (write) at (0,-2.5) {空闲结果槽（1/2 轮换）\\选择当前显示槽以外的槽};
\node[fpblock] (pending) at (5,-2.5) {待发布完整结果\\最后块末行写完成 → done};
\node[fpblock] (show) at (10,-2.5) {显示持有结果槽\\按独立光栅重复读出};
\draw[fparrow] (write)--(pending);
\draw[fparrow] (pending)--(show);
\draw[fparrow] (show.south)--++(0,-0.9)-| (write.south);
\node[fpnote,fill=white] at (5,-3.65) {新结果替换显示槽后，旧显示槽方可写入下一结果};
\node[fpnote,text width=14cm] at (5,-1.2)
 {发布门槛：显示 VS 请求新 epoch → 旧读事务结束 → 旧 FIFO 排空 → 确认与接受握手};
\end{tikzpicture}
\caption{源帧与结果帧的逻辑所有权流转（概念状态，与 RTL 标志及控制状态对应）}
\label{fig:fpga-frame-ownership}
\end{figure}


\subsection{结果发布与显示 epoch 握手}
结果帧始终写入当前显示槽之外的 1/2 号槽。
\path{renderer_done} 且无错误时设置 \path{pending_valid}，表示完整结果已提交。
显示域在 VS 上升沿翻转 \path{display_request}，开始刷新 epoch，停止使用旧像素并排空显示 FIFO。
DDR 域检测到新 epoch 后停止发旧帧读命令，等在途读完成且 FIFO 写域空闲额度恢复为全深度，
才确认 epoch；此时如有待发布结果则切换显示槽，否则继续使用原显示槽。
显示域收到确认后返回 \path{display_accept_epoch}，DDR 域收到此接受消息才允许新一轮预取。
因此旧事务、旧 FIFO 数据均先结束，随后才允许新的帧数据进入。

显示按最多 16 beat 突发预取，发命令前检查 FIFO 空闲额度能够容纳整笔返回。
若有效扫描期间 FIFO 饥饿，置位粘滞诊断 \path{display_underrun}，
并把本帧余下像素补黑、排掉迟到数据，下一帧边界重新交接，避免晚到像素造成水平错位。
\path{displayed_frames} 只在发布新结果时增加，不能用它直接计数全部物理 HDMI 扫描帧。
\sourcefile{\path{../Pgl50h/src/video/pose/fram_buf_pose.v}}

\section{双槽分块渲染器设计}
\subsection{接口契约、块划分与片内存储}
\path{pose_renderer_tiled} 工作在 100 MHz DDR 用户时钟域。
\path{start} 仅在 \path{busy=0} 时接收；同时锁存
\path{source_bank[1:0]}、\path{render_bank[1:0]}、
\path{yaw_q2048/pitch_q2048[15:0]} 和 \path{lens_enable}。
整个渲染过程中这些帧级参数保持不变。
\path{done} 为一拍完成脉冲，\path{completed_bank[1:0]} 指示结果槽；
\path{error} 表示容量或数据返回异常，出现后停止继续调度，当前实现依赖复位恢复。

输出划分为 $128\times32$ 的块，按从左到右、从上到下顺序遍历。
1080p 共 15 列、34 排、510 块，最后一排高度为 24。
每块所有输出像素均执行定点逆投影，逐点记录最近邻采样坐标，
不采用只计算块角点的网格近似。
内部块槽 0/1 各包含坐标 RAM、源区域 RAM 和输出 RAM，容量见表~\ref{tab:fpga-tile-memory}。
\begin{table}[htbp]\centering\small
\caption{默认分块配置的逻辑缓存容量}\label{tab:fpga-tile-memory}
\begin{tabularx}{\textwidth}{p{4.0cm}p{4.4cm}Y}
\toprule 缓存 & 容量 & 实现用途 \\\midrule
\path{coordinates} & 总计 $8192\times23$ bit & 每项为有效位、11 bit 源行坐标、11 bit 源列坐标。\\
\path{source_mem} & 总计 $4096\times256$ bit & 每槽 2048 beat，即 64 KiB 源矩形缓存。\\
\path{output_mem} & 总计 $512\times256$ bit & 每槽 256 beat，容纳一个完整 $128\times32$ 结果块。\\
采集/显示 FIFO & 各 $512\times256$ bit & 各 16 KiB，跨像素域与 DDR 域。\\
采集/渲染写 RAM & 各 $16\times256$ bit & 客户端持有的分布式 RAM 突发缓冲。\\
\bottomrule
\end{tabularx}
\end{table}
这些是 RTL 逻辑容量，实际 DRM 位置数还取决于端口组织和工具映射。
大容量 RAM 不逐项复位，而是复位指针、占用和有效状态，
以先写后读及槽所有权防止使用旧内容，并保持同步块 RAM 推断。

\subsection{投影、包围矩形与源区域预取}
图~\ref{fig:fpga-tile-flow} 将投影/读源、采样和写回三个并行状态机展开。
\path{pb} 指向投影/预取块槽，\path{sb} 指向采样/写回块槽；
\path{occupied} 与 \path{ready_tiles} 分别表示槽被占用、坐标和源区域已准备好。
投影器每拍接收一个块内输出位置，返回值有效时用
\begin{equation}
x_t=\left\lfloor u_s+\tfrac12\right\rfloor,\qquad
 y_t=\left\lfloor v_s+\tfrac12\right\rfloor
\end{equation}
得到最近邻坐标。RTL 直接使用 Q14 坐标整数部分加第 13 位实现舍入，
并对所有有效采样点求 $x_{\min},x_{\max},y_{\min},y_{\max}$。
\begin{figure}[htbp]
\centering
\begin{tikzpicture}[fpblock/.append style={text width=3.8cm,font=\fontsize{8.5}{11}\selectfont,minimum height=10mm}]
\node[fpnote] at (0,0.95) {投影 / 预取状态机 \texttt{ps}};
\node[fpnote] at (5,0.95) {采样状态机 \texttt{ss}};
\node[fpnote] at (10,0.95) {写回状态机 \texttt{ws}};
\node[fpblock] (slot) at (0,0) {准备姿态，等待空闲块槽\\\texttt{P\_POSE / P\_SLOT}};
\node[fpblock] (proj) at (0,-1.6) {逐点投影，坐标写 RAM\\累积有效采样点的包围矩形};
\node[fpblock] (area) at (0,-3.2) {\texttt{P\_BOUNDS → P\_AREA}\\\texttt{→ P\_CAPACITY}\\超过 2048 beat 则置错};
\node[fpblock] (read) at (0,-4.8) {按源矩形逐行读 DDR\\\texttt{P\_READ\_CMD / DATA}};
\node[fpblock] (ready) at (0,-6.4) {源区域完整 → 发布块就绪\\\texttt{P\_READY}，翻转 \texttt{pb}};
\node[fpblock] (wait) at (5,0) {\texttt{S\_WAIT}\\等待 \texttt{ready\_tiles[sb]}};
\node[fpblock] (coord) at (5,-1.6) {读坐标 RAM\\取得有效位、源 $x/y$};
\node[fpblock] (sample) at (5,-3.2) {相对地址 → 源 RAM\\按 lane 取样，无效点填黑};
\node[fpblock] (pack) at (5,-4.8) {16 像素打包写输出 RAM\\完整行递增 \texttt{rows\_ready}};
\node[fpblock] (drain) at (5,-6.4) {\texttt{S\_DRAIN}\\等待该块写完后翻转 \texttt{sb}};
\node[fpblock] (row) at (10,0) {\texttt{W\_WAIT}\\等待下一输出行就绪};
\node[fpblock] (load) at (10,-1.6) {\texttt{W\_LOAD}\\同步 RAM → 客户端行缓冲};
\node[fpblock] (cmd) at (10,-3.2) {\texttt{W\_CMD}\\提交写命令并保持行缓冲};
\node[fpblock] (done) at (10,-4.8) {\texttt{W\_DONE}\\等待真实末拍 \texttt{wr\_done}};
\node[fpblock] (release) at (10,-6.4) {行未结束则继续下一行\\块结束释放槽，帧结束置 done};
\foreach \a/\b in {slot/proj,proj/area,area/read,read/ready,wait/coord,coord/sample,sample/pack,pack/drain,row/load,load/cmd,cmd/done,done/release}
  \draw[fparrow] (\a)--(\b);
\draw[fparrow,dashed] (ready.east)--(2.5,-6.4)--(2.5,0)--(wait.west);
\node[fpnote,rotate=90,fill=white,font=\scriptsize] at (2.5,-2.2) {块就绪后消费};
\draw[fparrow,dashed] (pack.east)--(7.5,-4.8)--(7.5,0)--(row.west);
\node[fpnote,rotate=90,fill=white,font=\scriptsize] at (7.5,-2.2) {行就绪即可写回};
\node[fpnote,text width=14cm] at (5,-7.5) {全块无效时跳过源读取；两个块槽交替：下一块投影/预取与当前块采样/写回重叠。};
\end{tikzpicture}
\caption{双槽分块渲染流程（按三个实际状态机分组，省略等待自环）}
\label{fig:fpga-tile-flow}
\end{figure}


源矩形按 16 像素打包边界横向扩展，所需 beat 数为
\begin{equation}
W_s=\left\lfloor\frac{x_{\max}}{16}\right\rfloor-
\left\lfloor\frac{x_{\min}}{16}\right\rfloor+1,\qquad
H_s=y_{\max}-y_{\min}+1,\qquad B_s=W_sH_s.
\end{equation}
\path{P_BOUNDS} 寄存矩形宽高，\path{P_AREA} 寄存面积，
\path{P_CAPACITY} 检查 $B_s\le2048$，将减法、乘法和比较分到不同周期。
超限时在发出首个读命令前置错，避免源 RAM 地址截断引起静默覆盖。
全块坐标无效时直接标记块就绪，后续生成黑块而不读源帧。

读取接口为 \path{rd_cmd_valid/ready}、\path{rd_cmd_addr[27:0]}、
\path{rd_cmd_len[4:0]}，长度为实际 beat 数 1～16。
按源矩形逐行发突发，不把跨行的无关数据并入同笔请求。
返回接口 \path{rd_data_valid}、\path{rd_data[255:0]}、\path{rd_data_last}
没有回压；这里的 \path{rd_data_last} 接的是仲裁器核验后的读完成脉冲。
每拍写入当前块槽源 RAM，并核对末拍位置与累计缓存容量。
读源完成后发布 \path{ready_tiles[pb]}，投影端转向另一空槽。

\subsection{最近邻采样、行就绪与结果写回}
采样状态机对就绪块依序读取坐标，经过“坐标 RAM → 相对地址 → 源 RAM → lane 选择 → 打包”
流水线，每拍启动一个像素。若源矩形左边界的 beat 列为 $x_{w0}$、首行为 $y_0$，
局部源 RAM 地址为
\begin{equation}
a_{\mathrm{local}}=(y_t-y_0)W_s+
\left(\left\lfloor x_t/16\right\rfloor-x_{w0}\right).
\end{equation}
以 $x_t\bmod16$ 选择 16 bit 像素；坐标有效位为零时输出黑色。
每 16 个像素打包写入输出 RAM，\path{rows_ready} 记录已完整生成的行数。

写回状态机无需等待整块采样结束，只需下一行已经完整生成，即可在
\path{W_LOAD} 将该行从同步输出 RAM 装载到客户端突发 RAM，
排空最后一拍 RAM 读延迟后，以 \path{W_CMD} 发写请求。
一个 128 像素行对应 8 beat，写地址按式~\eqref{eq:fpga-ddr-address} 计算。
\path{wr_cmd_valid/ready} 仅表示仲裁器接收命令，
客户端仍须保持缓冲内容，直到 \path{wr_done} 才推进结果行。
块最后一行真实写完后清除占用、切换 \path{sb}；最后一个块提交后才撤销
\path{busy} 并产生 \path{done}，不会在仅完成投影或发出写地址时发布结果。

两个块槽允许第 $n+1$ 块投影及读源与第 $n$ 块采样和写回重叠；
同块内还允许已完成行写回与后续行采样重叠。
当两个几何开关均关闭、帧姿态为零时，渲染器选择精确恒等坐标通路，
避开定点往返误差和投影器的边界保护，保留最后一行、最后一列的直接拷贝。
\sourcefile{\path{../Pgl50h/src/video/pose/pose_renderer_tiled.v}}

\section{姿态与镜头联合逆投影流水线}
\subsection{模块接口与定点格式}
\path{pose_inverse_project_q14} 以 \path{pose_load} 锁存帧姿态及镜头开关，
准备三角系数后置 \path{pose_ready}；渲染器等待该信号后才发送像素。
输入为 \path{in_valid} 与 \path{dst_x/dst_y[10:0]}，
输出为 \path{out_valid}、\path{map_valid} 和两个有符号 32 bit Q14 源坐标。
\path{out_valid} 表示该拍对应一个已经处理完的输入坐标，
\path{map_valid} 才表示此坐标能够采样。无效位置仍占一个输出顺序位置，
下游将其填黑，不能删去该项；投影流水线没有 \path{ready} 回压接口。

\begin{table}[htbp]\centering\small
\caption{主要定点量及数值含义}\label{tab:fpga-fixed-point}
\begin{tabularx}{\textwidth}{p{4.0cm}p{3.2cm}Y}
\toprule 量 & 位宽/小数位 & 说明 \\\midrule
yaw、pitch & 有符号 16 bit，11 位小数 & Q2048 弧度值；实际角度为码值除以 2048。\\
正弦、余弦 & 有符号 16 bit，14 位小数 & 整数角度查表及表间线性插值。\\
归一化坐标、旋转射线 & 有符号 24 bit，18 位小数 & 乘法保留宽中间结果，按级移位回到 Q18。\\
$f_x,f_y,c_x,c_y$ & 有符号 32 bit，14 位小数 & 由标定参数预量化；焦距倒数为编译期常数。\\
$k_1,k_2,p_1,p_2$ & 有符号 18 bit，14 位小数 & 默认 $162,-711,0,0$。\\
输出源坐标 & 有符号 32 bit，14 位小数 & 下游按整数部分及最高小数位进行最近邻舍入。\\
\bottomrule
\end{tabularx}
\end{table}
默认渲染器向投影器传入 \path{POSE_LIMIT_Q2048=536}，即约
$\pm15^\circ$；超限输入先钳位并置 \path{pose_clamped}。
该上限来自调用实例，不能误用投影模块独立声明中更大的兼容默认值。
roll 虽由 IMU 解码，但默认联合矫正只使用 yaw、pitch。

\subsection{分级运算与有效性传播}
图~\ref{fig:fpga-projector} 给出联合映射的硬件运算顺序。
输出坐标从 0 开始，而标定主点采用原有的一基坐标，因此归一化使用
$(u_d+1-c_x)/f_x$、$(v_d+1-c_y)/f_y$，投回源图时再减 1。
归一化乘以编译期预求的焦距倒数，不在像素通路中计算焦距除法。
\begin{figure}[htbp]
\centering
\begin{tikzpicture}
\node[fpblock] (norm) at (0,0) {输出像素归一化\\一基 / 零基换算、常数乘法};
\node[fpblock] (rotate) at (5,0) {yaw / pitch 射线旋转\\前端共五级寄存};
\node[fpblock] (divide) at (10,0) {透视归一化 $x/z,y/z$\\两个并行 42 级恢复除法器};
\node[fpblock] (lens) at (10,-2.2) {径向与切向镜头模型\\$r^2,r^4$、系数乘积及求和};
\node[fpblock] (focal) at (5,-2.2) {焦距乘法与主点累加\\部分积拆分、逐级寄存};
\node[fpblock] (check) at (0,-2.2) {坐标范围 / 有效性判定\\Q14 坐标 + valid + map};
\node[fpnote] (pose) at (5,1.35) {帧级 \texttt{pose\_load} → 钳位 / 查表插值 → \texttt{pose\_ready}};
\draw[fparrow,dashed] (pose)--(rotate);
\draw[fparrow] (norm)--(rotate);
\draw[fparrow] (rotate)--(divide);
\draw[fparrow] (divide)--(lens);
\draw[fparrow] (lens)--(focal);
\draw[fparrow] (focal)--(check);
\node[fpnote,text width=14cm] at (5,-3.4) {valid 与候选有效位随各级传播；流水线每拍可接收一个坐标，无效映射仍保留输出位置。};
\end{tikzpicture}
\caption{姿态与镜头联合逆投影的定点硬件流水线}
\label{fig:fpga-projector}
\end{figure}


三角系数准备按“幅值钳位 → 弧度转角度 → 相邻整数角度查表 → 插值乘积 → 符号恢复”分级。
像素前端分成归一化、yaw 乘积、yaw 求和、pitch 乘积、pitch 求和五级，
实现第三章所定义的逆向射线映射
$\boldsymbol{r}=R_x(\mathrm{pitch})R_y(\mathrm{yaw})[x_d,y_d,1]^\mathsf{T}$。
这里采用从理想输出射线映射到当前相机射线的方向，
不将用于描述相机朝向的逆矩阵 $C=R_y(-\mathrm{yaw})R_x(-\mathrm{pitch})$ 直接用于采样。
系数快照随像素前端一起寄存，避免控制更新与在途像素混用。

透视归一化使用两个并行的恢复除法器计算 $r_x/r_z$、$r_y/r_z$。
分子宽度为 $24+18=42$ bit，每一级处理一位，形成 42 级除法流水线，
从而用多级延迟换取每拍一组坐标的吞吐，而非每像素等待一次串行除法。
默认关闭实验性的倒数 LUT 分支。
射线须满足 $r_z>1024/2^{18}$ 且 $|r_x|,|r_y|<2r_z$，
不满足时仍传播有效拍但清除候选有效位，并使用安全分母/分子进入后续运算。

镜头模型依次计算 $x^2,y^2,xy,r^2,r^4$、径向系数和切向偏移，
再乘焦距并加主点。乘法、交叉项求和与主点累加分别寄存，
较宽焦距乘法还拆分为高低部分积，缩短单周期组合链。
镜头开关关闭时径向缩放设为 1、切向项设为 0，姿态映射仍正常工作。
镜头开启时还要求 $r^2\le600000/2^{18}$。
最终保守有效范围为 $0\le u_s<1919$、$0\le v_s<1079$，
这是投影模块沿用的四邻域边界保护；默认最近邻路径也继承该判据，
仅前述精确恒等通路绕过它。所有阶段同步传递 valid 和候选标志，
渲染器以返回 valid 计数排空流水线，不用固定等待估计完成时间。
\sourcefile{\path{../Pgl50h/src/video/pose/pose_inverse_project_q14.v}}

\section{DDR3 仲裁与事务完成控制}
\subsection{用户接口与客户端握手}
\path{pose_ddr_arbiter} 为采集写、渲染写、显示读、渲染读提供统一接口。
写通道固定采集优先，读通道固定显示优先；读写通道分别拥有独立状态机，
允许同时存在一笔读事务与一笔写事务，各自最多一笔在途，不交错同向事务。
同向优先请求持续存在时，渲染请求等待仲裁。

\begin{table}[htbp]\centering\small
\caption{DDR3 IP 用户接口的实际语义}\label{tab:fpga-ddr-interface}
\begin{tabularx}{\textwidth}{p{5.0cm}p{2.2cm}Y}
\toprule 信号 & 方向（仲裁器） & 语义 \\\midrule
\mbox{\texttt{axi\_awaddr[27:0]}}、\mbox{\texttt{awid[3:0]}}、\mbox{\texttt{awlen[3:0]}}、\mbox{\texttt{awvalid}} & 输出 & 写首地址、客户端 ID、beat 数减 1、地址有效；等待 \path{awready}。\\
\mbox{\texttt{axi\_wdata[255:0]}}、\mbox{\texttt{wstrb[31:0]}} & 输出 & 当前 beat 及字节写使能，默认屏蔽全开。\\
\mbox{\texttt{axi\_wready}}、\mbox{\texttt{axi\_wlast}} & 输入 & IP 接受当前写 beat，IP 指示该拍为最后一拍。\\
\mbox{\texttt{axi\_araddr[27:0]}}、\mbox{\texttt{arid[3:0]}}、\mbox{\texttt{arlen[3:0]}}、\mbox{\texttt{arvalid}} & 输出 & 读首地址、客户端 ID、beat 数减 1、地址有效；等待 \path{arready}。\\
\mbox{\texttt{axi\_rdata[255:0]}}、\mbox{\texttt{rid[3:0]}}、\mbox{\texttt{rvalid/rlast}} & 输入 & 无回压读返回，逐 beat 核对 ID 与最后一拍标志。\\
\bottomrule
\end{tabularx}
\end{table}
该接口是厂商 DDR3 IP 的精简 AXI 风格接口，不能按完整 AXI4 连接理解：
没有 \path{WVALID}、\path{RREADY} 或写响应 B 通道，
\path{axi_wlast} 是 IP 输出给仲裁器的信号。
因此发 AW 前客户端必须准备好整笔写数据，发 AR 前必须预留整笔返回空间。
内部命令长度为 5 bit 的实际 beat 数 $L\in[1,16]$，
对 IP 编码为 4 bit 的 $L-1$；采集/显示 ID 为 0，渲染 ID 为 1。

\subsection{事务状态机及共享写缓冲}
如图~\ref{fig:fpga-ddr-fsm}，\path{IDLE} 接受客户端
\path{valid && ready} 后锁存地址、长度和所有者；\path{ADDRESS} 保持地址及 valid，
直至 IP 接受；\path{PAYLOAD} 按真实 beat 传输推进计数。
实现还覆盖地址握手与首个数据传输出现在同一周期的情况。
\begin{figure}[htbp]
\centering
\begin{tikzpicture}[fpstate/.append style={text width=4cm}]
\node[fpstate] (idle) at (0,0) {\texttt{IDLE}\\客户端握手，锁存合法命令};
\node[fpstate] (address) at (5,0) {\texttt{ADDRESS}\\地址、ID、长度不变\\VALID 保持至握手};
\node[fpstate] (payload) at (10,0) {\texttt{PAYLOAD}\\逐 beat 核验并推进计数};
\node[fpstate] (fault) at (10,-2.1) {\texttt{FAULT}\\锁存错误，停止本通道};
\draw[fparrow] (idle)--(address);
\draw[fparrow] (address)--(payload);
\draw[fparrow] (payload.north)--++(0,1.0)-| (idle.north);
\node[fpnote,fill=white] at (5,1.25) {正确末拍：产生客户端 done，返回空闲};
\draw[fparrow] (payload)--node[right,font=\footnotesize]{异常 LAST / ID}(fault);
\draw[fparrow] (idle.south)--(0,-2.1)--node[above,font=\footnotesize]{非法长度；异常返回同样进入故障}(fault.west);
\node[fpnote,text width=14cm] at (5,-3.45) {写：\texttt{WREADY} 接受数据，核验 IP 返回的 \texttt{WLAST}；读：\texttt{RVALID} 时接收并核验 \texttt{RID/RLAST}。\\地址接受与首拍数据可以同周期发生；故障通道只通过复位恢复。};
\end{tikzpicture}
\caption{DDR 事务状态与完成条件（读写通道分别独立运行）}
\label{fig:fpga-ddr-fsm}
\end{figure}


每次 \path{WREADY} 有效才推进写 beat 索引，并核对
\path{WLAST==(remaining==1)}；只有正确接受最后一个 beat 才产生
\path{cap_w_done} 或 \path{render_w_done}。
此处完成的含义是 IP 用户接口已接受全部数据，并非 DDR 器件物理写入的独立确认。
读返回则要求当前处于允许返回的事务阶段、ID 匹配且 LAST 与剩余长度一致，
通过后路由到对应客户端，末拍同时产生 \path{*_r_done}。
非法长度、过早/过晚 LAST、错误 ID 或无事务返回均使相应通道进入
\path{FAULT}，置位 \path{protocol_error}，直到复位才恢复，防止错误事务引起帧槽提前复用。

默认 \path{SHARED_WRITE_PAYLOAD=1}，采集端与渲染端各自持有
$16\times256$ bit 分布式 RAM，仲裁器输出 \path{write_beat_index[3:0]}，
从当前所有者读出 \path{cap_w_beat/render_w_beat[255:0]}。
等待 \path{WREADY} 时索引不变，数据保持稳定；客户端直到 done 才重填缓冲。
旧式 \path{wr_data[4095:0]} 端口在该模式固定为零，
避免重复存储整笔 4096 bit 载荷，降低宽寄存器和选择网络开销。
\sourcefile{\path{../Pgl50h/src/video/pose/pose_ddr_arbiter.v}；\path{fram_buf_pose.v}；\path{pose_renderer_tiled.v}}

\section{姿态输入、运行控制与统计遥测}
\subsection{IMU 协议及帧姿态快照}
\path{imu_uart_rx_50m} 在 50 MHz 域同步异步串口输入并按 16 倍波特率采样。
\path{imu_protocol_50m} 依次识别
\texttt{55 AA}、6 字节姿态载荷、载荷累加和低 8 bit、帧尾 \texttt{0D}。
yaw、pitch、roll 各占小端序 16 bit；仅校验和及帧尾均正确时，
同时更新三轴数据、产生一拍 \path{attitude_valid} 并翻转 \path{attitude_toggle}。

\path{attitude_cdc} 将事件翻转两级同步并捕获稳定角度快照。
顶层同时保留向输出像素域和输入像素域的同步实例；默认完整帧路径使用输入域的快照，
在被选中源帧的 SOF 锁存 yaw/pitch，此后直到源帧提交与渲染调度均保持不变。
姿态开关关闭时顶层向帧缓存送零角度；实现使用绝对姿态值，未在该通路额外加入三轴融合、
roll 补偿或相邻帧姿态差计算。
该 CDC 没有返回确认，正确性依赖 IMU 更新间隔足够长、数据在快照捕获期间稳定，
不能用于任意高速变化总线。

\subsection{PC 控制的帧边界原子应用}
控制掩码 \path{mask[2:0]} 中 bit 0、1、2 分别为镜头、姿态、自动 Gamma 使能。
串口包为 \texttt{A5 5A | version | type | sequence | length | payload | XOR}，
版本为 1，XOR 覆盖 version 到 payload。
\path{cp2102_control} 校验长度、可用功能、参数范围和忙状态后，
输出 \path{control_valid} 与 \path{control_mask[2:0]}。
\begin{table}[htbp]\centering\small
\caption{PC 串口命令与应答类型}
\begin{tabularx}{\textwidth}{p{2.2cm}p{3.0cm}Y}
\toprule 类型 & 载荷 & 行为 \\\midrule
\texttt{0x10} & 1 字节功能掩码 & 请求更新三个功能开关，忙时拒绝重叠请求。\\
\texttt{0x11} & 无 & 查询请求值、应用值及控制状态。\\
\texttt{0x12} & 1 字节 0/1 & 关闭/开启周期遥测。\\
\texttt{0x13} & 无 & 请求最近一次统计快照，无快照时返回错误。\\
\texttt{0x80} & 8 字节 & 状态应答，保留请求序号并携带结果码。\\
\texttt{0x81} & 28 字节 & 统计遥测，使用独立递增序号。\\
\bottomrule
\end{tabularx}
\end{table}

\path{video_control_cdc} 在系统域锁存 \path{requested_mask} 并翻转 request，
输入域同步后在寄存 VS 产生的 SOF 才更新 \path{applied_mask_vin} 并返回 acknowledge；
系统域收到确认后才更新 \path{applied_mask_sys} 并清 busy。
图~\ref{fig:fpga-control-cdc} 给出顺序。
Gamma 在该 SOF 使用 \path{next_mask_vin}，确保与掩码提交同帧生效；
后级几何参数在延迟后的 RGB565 SOF 使用已应用掩码。
因此同一源帧的三个开关一致，控制值保持到整帧结束。
\begin{figure}[htbp]
\centering
\begin{tikzpicture}[fpblock/.append style={text width=4.9cm}]
\node[fpnote] at (0,0.9) {50 MHz 系统域};
\node[fpnote] at (8,0.9) {输入像素域};
\node[fpblock] (req) at (0,0) {锁存 \texttt{requested\_mask[2:0]}\\翻转 request，置 busy};
\node[fpblock] (sync) at (8,0) {request 两级同步\\检测到尚未应用的新请求};
\node[fpblock] (hold) at (0,-1.8) {保持掩码直至确认返回\\忙时拒绝后续重叠命令};
\node[fpblock] (wait) at (8,-1.8) {等待寄存输入 VS 的 SOF\\帧内不改变应用掩码};
\node[fpblock] (ack) at (0,-3.6) {acknowledge 两级同步\\更新系统域应用值，清 busy};
\node[fpblock] (apply) at (8,-3.6) {原子更新 \texttt{applied\_mask\_vin}\\返回 acknowledge 翻转};
\draw[fparrow] (req)--node[above,font=\footnotesize]{稳定数据 + 事件}(sync);
\draw[fparrow] (req)--(hold);
\draw[fparrow] (hold)--(ack);
\draw[fparrow] (sync)--(wait);
\draw[fparrow] (wait)--(apply);
\draw[fparrow] (apply)--node[above,font=\footnotesize]{应用后确认}(ack);
\node[fpnote,text width=14cm] at (4,-4.8) {Gamma 通过 \texttt{next\_mask\_vin} 在提交 SOF 取新值；几何开关与姿态在后级源帧 SOF 锁存。};
\end{tikzpicture}
\caption{PC 控制掩码跨域握手与帧边界应用流程}
\label{fig:fpga-control-cdc}
\end{figure}


命令应答成功表示请求已接收，是否已应用需结合 busy 和应用掩码判断；
LED 显示源帧侧的已应用状态，显示器仍可能正在扫描较早配置生成的完整结果帧。
输入无新 SOF 时请求保持 busy，不会在帧中途强制应用。

\subsection{遥测快照与跨域发送}
\path{gamma_telemetry} 同时接收原始统计和 Gamma 输出像素，
额外统计输出亮度和；亮度计算先寄存再累加，将 ROM 查询与 32 bit 累加链隔开。
每份 224 bit 快照包含帧号、输入亮度和、有效像素数、暗/高亮像素数、
输出亮度和，以及各占一字节的曲线号、目标号、使能和掩码。
统计元数据对齐已完成帧，避免把新帧开关附在旧帧亮度上。

\path{cdc_snapshot} 使用数据保持寄存器和 request/acknowledge 翻转握手，
将整份快照从输入像素域送到 50 MHz 域；源端 busy 时丢弃新采样，保持在途数据完整。
\path{cp2102_control} 缓存最近快照，默认 10 Hz 上报，命令应答优先于遥测。
多字节字段按低字节先发。这条通路用于观察亮度和开关，
不包含渲染完成率、DDR 错误计数或实板帧率测量。
\sourcefile{\path{../Pgl50h/src/imu/}；\path{../Pgl50h/src/control/}}

\section{跨时钟域实现与工程约束}
\subsection{异步 FIFO 的存储与指针同步}
采集与显示各使用一份 \path{pose_async_fifo}，默认宽 256 bit、深 512，
二进制读写指针均为 10 bit（含回绕位），跨域前转换为 Gray 码：
\begin{equation}
g=b\oplus(b\gg1).
\end{equation}
远端 Gray 指针经两级同步后用于空满判断及本域额度计算。
读端使用同步 RAM 预取当前或下一队首，只有
\path{rd_valid && rd_ready} 才前移读指针；停顿时保持队首，
预取字仍计入占用，不会被写端提前覆盖。
\path{wr_free} 为写域保守空闲额度，\path{rd_count} 为读域保守占用计数，
均有跨域同步延迟，不能当作另一时钟域的瞬时状态。

两域共用异步复位断言，各自同步释放；复位丢弃所有 FIFO 数据，
RAM 内容不清零，只有 valid 和指针决定内容能否读取。
除 FIFO 外，配置和遥测采用有确认邮箱，IMU 采用事件翻转快照，
采集完成及显示 epoch 使用专门握手。时序例外只描述时钟关系，不能替代这些同步结构。

\subsection{FDC 与实现检查的对应关系}
\path{src/constraints/hdmi_ddr3_top.fdc} 为顶层分配引脚、电平标准、
输入/输出寄存器位置和时钟约束，主要设计关系见表~\ref{tab:fpga-fdc}。
\begin{table}[htbp]\centering\small
\caption{RTL 结构与工程时序约束的对应关系}\label{tab:fpga-fdc}
\begin{tabularx}{\textwidth}{p{3.5cm}p{4.3cm}Y}
\toprule 对象 & 工程约束 & 对应实现含义 \\\midrule
主/生成时钟 & 50 MHz、27 MHz、输入周期 6.734 ns；声明视频和 DDR 生成时钟 & 保留 PLL 倍分频及 DDR PHY 时钟来源。\\
HDMI 输入 & 输入延迟最大 2.000 ns、最小 $-0.500$ ns & 相对于 \path{pixclk_in} 的源同步预算。\\
输入首级捕获 & 对 RGB、DE、VS 外部到首级路径设置 setup multicycle 0 & 描述同一发射周期的 I/O 寄存捕获；后续内部路径仍为普通单周期。\\
HDMI 输出 & 输出延迟最大 1.500 ns、最小 $-0.500$ ns；约束时钟反相 & 用反相生成时钟描述 MS7210 下降沿采样，物理输出时钟保持正相。\\
异步时钟组 & 按板载、输入、视频、DDR PHY 关系划组 & 每条跨域路径对应 FIFO、同步器或厂商 IP 结构。\\
FIFO Gray 总线 & 四条总线各 10 bit，最大偏斜 6.734 ns & 限制指针位间到达偏斜，独立于异步时钟组审查。\\
HDMI 输出 I/O & \path{PAP_IO_REGISTER=TRUE} & 缩短 fabric 至管脚路径，并控制 RGB/同步总线偏斜。\\
\bottomrule
\end{tabularx}
\end{table}
输入同周期捕获依赖当前板级输入时序预算与时钟布线关系，
更换桥芯片配置、板卡或输入频率时需重新核对该假设。
PDS 工程通过 \path{ipcore/} 下的 \path{.idf} 引入 PLL 和 DDR3 IP，
因此重建工程还需保留 IP 配置与 FDC，不能仅复制应用 RTL。

现有 2026 年 10 月 22 日实现归档包含快慢工作角时序及四条 Gray 指针偏斜审计，
本章用其说明当前代码与约束的组织方式；资源、处理帧率及仿真证据在后续章节分别讨论。
这些工具结果不代替板卡下载后的端到端验证，当前文档不将其表述为实板测量。
\sourcefile{\path{../Pgl50h/src/video/pose/pose_async_fifo.v}；\path{../Pgl50h/src/constraints/hdmi_ddr3_top.fdc}；\path{../Pgl50h/verification/20261022_pc_control_sbit.md}}
